\centerline{\bf \Huge Алгебраический разнобой}

\begin{flushright}
        \scriptsize{— Ты совсем не изменился.\\— Ну, конечно, я изменился. Ведь жить — значит меняться.\\Мисато и Кадзи}
\end{flushright}

\problem{1}
На координатной плоскости нарисована парабола $y=x^2$. Для данного числа $k>0$ рассматриваются трапеции, вписанные в эту параболу (то есть все вершины трапеции лежат на параболе), у которых основания параллельны оси абсцисс, а произведение длин оснований равно $k$. Докажите, что диагонали всех таких трапеций проходят через одну точку.

\problem{2}
Пусть $x_1<x_2<\ldots<x_{2026}$ - возрастающая последовательность натуральных чисел. При $i=1,2, \ldots, 2026$ обозначим $p_i=\left(x_1-\frac{1}{x_1}\right)\left(x_2-\frac{1}{x_2}\right) \ldots\left(x_i-\frac{1}{x_i}\right)$. Какое наибольшее количество натуральных чисел может содержаться среди чисел $p_1, p_2, \ldots, p_{2026}?$

\problem{3}
Дан квадратный трехчлен $P(x)$. Докажите, что существуют попарно различные числа $a, b$ и $c$ такие, что выполняются равенства

$$
P(b+c)=P(a), \quad P(c+a)=P(b), \quad P(a+b)=P(c) .
$$

\problem{4}
Множество $A$ состоит из $n$ различных натуральных чисел, сумма которых равна $n^2$. Множество $B$ также состоит из $n$ различных натуральных чисел, сумма которых равна $n^2$. Докажите, что найдётся число, которое принадлежит как множеству $A$, так и множеству $B$.

\problem{5}
Яна и Саша стартуют по круговой дорожке из одной точки в направлении против часовой стрелки. Обе бегут с постоянными скоростями, скорость Саши на $2\%$ больше скорости Яны. Яна всё время бежит против часовой стрелки, а Саша может менять направление бега в любой момент, непосредственно перед которым она пробежала полкруга или больше в одном направлении. Покажите, что пока Яна бежит первый круг, Саша может трижды, не считая момента старта, поравняться (встретиться или догнать) с ней.